% Corrigé intégré automatiquement pour la banque Exercices.
% Source : A_Integrer/DDS_04/071_PFS/071_PFS.tex
\begin{corrigebox}[Corrigé]
\CorrigeQuestion{1}
$P_B = \rho_v \pi \left(R^2 - R_M^2 \right)Lg$.
\CorrigeQuestion{2}
$P_B = \SI{101764}{N}$.
\CorrigeQuestion{3}
$\vect{F}+\vect{R_A}+\vect{R_B}+\vect{P_B}+\vect{T}=\vect{0}$.
\CorrigeQuestion{4}
On isole le solide ou l'ensemble indiqué, on dresse le bilan complet des actions mécaniques extérieures et on choisit un point de réduction qui élimine le maximum d'inconnues. Le PFS donne $\sum\vec F=\vec0$ et $\sum\vec M_A=\vec0$ ; les projections utiles fournissent les inconnues, puis leur signe permet de vérifier le sens supposé.
\CorrigeQuestion{5}
$C_m = rR_B$.
\CorrigeQuestion{6}
$T =\left(R_B  -R_A\right) \dfrac{\sin \left(\theta+\varphi\right)}{\cos\theta } $.
%$T = \dfrac{R_B \sin \left( \varphi+ \theta \right)- R_A \sin\theta }{\cos \theta }$
\CorrigeQuestion{7}
On projette l'expression de la question 4 sur $\vect{y}$ et on a 

\noindent $R_A \cos \left(\theta+\varphi\right) - R_B\cos \left(\theta+\varphi\right) - T \sin\theta -F(t)-P= 0$ 

soit

$F(t)=R_A \cos \left(\theta+\varphi\right) - R_B\cos \left(\theta+\varphi\right) - T \sin\theta -P$.

En utilisant le résultat de la question précédente, on a : 

$F(t)=R_A \cos \left(\theta+\varphi\right) - R_B\cos \left(\theta+\varphi\right) - \left(R_B  -R_A\right) \dfrac{\sin \left(\theta+\varphi\right)}{\cos\theta } \sin\theta -P$.

et

$F(t)=R_A \cos \left(\theta+\varphi\right) - R_B\cos \left(\theta+\varphi\right) - \left(R_B  -R_A\right)\sin \left(\theta+\varphi\right)  \tan \theta -P$.
\end{corrigebox}
